\section{Planning and Shared Objects}
\label{sec:planning}

\subsection{Task generation and plan finalization}

Service derives preparation tasks from the inferred household need and the
current household capability snapshot. Task generation identifies the need,
maps it to supported logical device actions, assigns trigger/timing and an
end condition, applies relevant household memory, runs deterministic
planning/policy guards, and finalizes versioned tasks. Unsupported or
product-prohibited tasks are removed before plan finalization.
\sourcecite{draft}

The baseline constrains generated actions through supported capabilities and
guards on values, timing, and rooms. Money-spending and irreversible actions
fall outside the intended autonomous-action scope. \sourcecite{baseline}

\subsection{Plan and task meanings}

A plan has an identity and version. A newer plan version supersedes older
execution intent. Service and Home must share the following task meanings.
\sourcecite{draft}

{\small
\begin{tabularx}{\textwidth}{@{}L{0.25\textwidth}Y@{}}
\toprule
\textbf{Meaning} & \textbf{Purpose} \\
\midrule
Plan identity & Associates the task with its plan or episode. \\
Plan version & Identifies the revision and supports stale-execution prevention. \\
Task identity & Correlates execution, cancellation, restoration, and results. \\
Target & Identifies the logical household device or capability. \\
Desired effect & States the requested household state or effect. \\
Trigger & Specifies when the task becomes eligible. \\
Deadline / validity & Bounds when a new execution can start. \\
Preconditions & States physical conditions required before execution. \\
End / restore policy & Determines how a temporary effect terminates or restores. \\
\bottomrule
\end{tabularx}
\par
}

The table specifies interoperable meanings. Exact JSON names, enum spellings,
and transport bindings remain open. \sourcecite{draft}

\subsection{Capability, current state, and product policy}

\proposed{Separate the current catalog into connector capability, current
Home state, and Service/shared product policy. Planning consumes a
Home-provided capability snapshot. Home revalidates current capability and
state when execution becomes eligible.}

{\small
\begin{tabularx}{\textwidth}{@{}L{0.25\textwidth}Y@{}}
\toprule
\textbf{Boundary} & \textbf{Meaning} \\
\midrule
Home / connector capability & Supported logical commands, value ranges and
units, readable state, raw events, verification support, and manual-change
attribution support. \\
Home runtime state & Current reachability, observed state, observable use or
conflict, provenance, and observation freshness. \\
Service/shared policy & Permitted actions, timing/context restrictions,
quiet/noisy treatment, user-visible preparation, and protected operations. \\
\bottomrule
\end{tabularx}
\par
}

Within this proposal, unsupported capability means the connector cannot
perform an operation. Product-prohibited use means policy excludes an otherwise
possible operation. Temporary unavailability means a normally supported
operation cannot execute now. These distinctions preserve the reason a task
cannot proceed. Exact status names remain open. \sourcecite{draft}\sourcecite{capability}
